% Finite instantiation of the prior quotient/discriminator mechanism.
\subsection{A prescribed-domain discriminator calculation}
\label{subsec:fixed-ntt-discriminator}

The selected-domain results do not imply that the prescribed NTT domain
has no lower-threshold proximity-gap example.  The quotient and
discriminator argument of~\cite[Appendix~A, Lemma~3 and
Proposition~4]{kkh} supplies such an example directly.  The following
finite calculation uses the same mechanism, refining its collision
bound by cancelling shared roots.

Let $p=2130706433$, $E=\mathbb F_{p^6}$, and
$D=\mu_{262144}\subset\mathbb F_p^*$.  The number of elements of
$E$ having degree exactly six over $\mathbb F_p$ is
\[
 Q=p^6-p^3-p^2+p.
\]
For a power-of-two fiber size $m\mid131072$, with $r\le s$ below, put
\[
 s=262144/m-1,\quad r=131072/m+1,\quad w=m-1,
 \quad G=\mu_{s+1}\setminus\{1\},\quad L=\binom sr.
\]
Write $V_S(Z)=\prod_{a\in S}(Z-a)$ for an $r$-subset $S$ of $G$.
If $|S\setminus T|=|T\setminus S|=u>0$, cancelling their common
factors leaves a nonzero difference of monic degree-$u$ polynomials,
of degree at most $u-1$.  A full-degree pole $b$ cannot be a shared
base-field root.  Moreover, its Frobenius orbit has six elements.
Thus $V_S(b)=V_T(b)$ holds for at most
$6\lfloor(u-1)/6\rfloor$ full-degree poles.  Define
\[
 W=\sum_{u\ge1}\binom ru\binom{s-r}u
                      6\left\lfloor\frac{u-1}{6}\right\rfloor.
\]
Averaging the ordered collision energy over the $Q$ eligible poles
gives a pole with energy at most $L+LW/Q$.  Cauchy--Schwarz then
gives at least
\[
 \left\lceil\frac{LQ}{Q+W}\right\rceil
\]
distinct nonzero values $V_S(b)$.  If $LW<Q$, some pole has no
off-diagonal collisions, so all $L$ labels are distinct.

Set $Y=X^m$ and let $R$ be the monic locator of any $w$ points in
the reserved fiber $Y^{-1}(1)$.  The one-pole compiler
\[
 f=RY^{r-1},\qquad g=\frac R{Y-b},\qquad
 h_S=R\frac{(Y-b)Y^{r-1}+V_S(b)-V_S(Y)}{Y-b}
\]
has witnesses of degree at most $131071$ and residual
$RV_S(Y)/(Y-b)$.  Its exact source and common agreement is
$131072+m-1$, and each represented label has exact agreement
$131072+2m-1$.  These equalities follow by the root-count and
interpolation argument in the one-pole compiler.  The domain remains
the prescribed subgroup $D$ throughout; only the pole is selected.

For $m=1024$, the image bound is
\[
 1553276470857367347050498710657094691569048149670139791
\]
at exact nearby agreement $133119$.  More relevant to the desired
label count, $m=4096$ gives $s=63$, $r=33$, and $LW<Q$, hence
\[
 L=\binom{63}{33}=860778005594247069
\]
distinct labels, exact source/common agreement $135167$, and exact
nearby agreement $139263$.  At most one label is lost on converting
to non-endpoint affine mixtures.  The remaining count exceeds
$\lceil p^6/2^{128}\rceil=274980728111395088$.
The agreement is nevertheless $519$ short of $139782$.

This shortfall cannot be removed inside the fixed-cardinality one-pole
full-fiber compiler merely by increasing $m$: at $m\ge8192$ there
are at most $32$ fibers even without reserving one, and hence at most
$\binom{32}{16}=601080390$ subsets of any fixed cardinality.
This is far below the required label count.  The calculation is a
finite prescribed-domain application of the prior mechanism, not a
benchmark bound or a new asymptotic construction.
Its integer receipt is
\nolinkurl{research/norm_quotient_coverage/verify_fixed_ntt_discriminator.py}.
